\documentclass[12pt]{article}
\usepackage[utf8]{inputenc}
\usepackage{geometry}
\usepackage{listings}
\usepackage{xcolor}
\usepackage{booktabs}
\usepackage{float}
\usepackage{hyperref}

\geometry{a4paper, margin=1in}

\title{Scanner Implementation Comparison: Manual vs JFlex}
\author{Compiler Construction Group}
\date{\today}

\lstset{
    basicstyle=\ttfamily\small,
    breaklines=true,
    frame=single,
    keywordstyle=\color{blue},
    commentstyle=\color{gray},
    stringstyle=\color{green!50!black},
    numbers=left,
    numberstyle=\tiny,
    stepnumber=1
}

\begin{document}

\maketitle

\section{Overview}
This document presents a comparison between two lexical analyzer implementations for the \textbf{CustomLang} programming language:
\begin{enumerate}
    \item \textbf{Manual Scanner}: A hand-coded Java implementation using DFA logic.
    \item \textbf{JFlex Scanner}: A generated scanner using the JFlex tool and regex specifications.
\end{enumerate}

Both implementations were tested against a suite of 5 test files covering valid tokens, complex expressions, edge cases, partial implementations, and error handling.

\section{Implementation Differences}

\begin{table}[H]
    \centering
    \begin{tabular}{|p{4cm}|p{5cm}|p{5cm}|}
        \hline
        \textbf{Feature} & \textbf{Manual Scanner} & \textbf{JFlex Scanner} \\
        \hline
        \textbf{Architecture} & Custom State Machine (DFA) & Table-driven Lexer (Generated) \\
        \hline
        \textbf{Token Recognition} & Explicit \texttt{if-else} and \texttt{while} loops & Regular Expressions in \texttt{.flex} file \\
        \hline
        \textbf{Priority} & Implemented via code execution order & Implemented via rule order in spec \\
        \hline
        \textbf{Error Detection} & Custom checks (e.g., \texttt{scanNumber} validation) & Catch-all rules and specific error regexes \\
        \hline
    \end{tabular}
    \caption{Functional Comparison}
\end{table}

\section{Output Comparison \& Analysis}

Both scanners produce identical token streams for valid inputs. However, subtle differences exist in error reporting strategies, particularly for boundary cases.

\subsection{Case 1: Invalid Identifiers (e.g., "2Count")}
\textbf{Input:} \texttt{2Count}

\textbf{Manual Scanner Behavior:}
The manual scanner follows a strict "longest match" for specific types. It consumes \texttt{2} as an \texttt{INTEGER\_LITERAL}, stops at \texttt{C}, and then consumes \texttt{Count} as a separate \texttt{IDENTIFIER}.
\begin{itemize}
    \item Output: \texttt{<INTEGER\_LITERAL, "2">}, \texttt{<IDENTIFIER, "Count">}
\end{itemize}

\textbf{JFlex Scanner Behavior:}
The JFlex specification includes a dedicated error rule for invalid identifiers starting with digits. It matches the entire \texttt{2Count} sequence as an error.
\begin{itemize}
    \item Output: \texttt{ERROR: Invalid Identifier (must start with uppercase)}
\end{itemize}

\textbf{Conclusion:} The JFlex approach provides more semantic error reporting for this specific case, while the Manual approach adheres strictly to token separation.

\subsection{Case 2: Malformed Floats (e.g., "1.2.3")}
\textbf{Input:} \texttt{1.2.3}

\textbf{Manual Scanner Behavior:}
Detects the second decimal point during the number scanning phase and immediately reports an error for the specific lexeme.
\begin{itemize}
    \item Output: \texttt{ERROR: Malformed Literal: Multiple decimal points}
\end{itemize}

\textbf{JFlex Scanner Behavior:}
Depending on rule definitions, JFlex may match \texttt{1.2} as a valid float and then \texttt{.3} as a sequence of invalid characters (or a float missing an integer part), triggering the corresponding error rule.
\begin{itemize}
    \item Output: \texttt{FLOAT(1.2)}, \texttt{ERROR(.3)}
\end{itemize}

\section{Side-by-Side Output Demonstration}

The following table demonstrates the identical token stream produced by both scanners for the input: \texttt{Count = 42;} (Note: \texttt{=} is invalid in this version).

\begin{table}[H]
    \centering
    \begin{tabular}{|p{7cm}|p{7cm}|}
        \hline
        \textbf{Manual Scanner Output} & \textbf{JFlex Scanner Output} \\
        \hline
        \begin{verbatim}
<IDENTIFIER, "Count", Line: 1, Col: 1>
[Error] Line: 1, Col: 7, Lexeme: "="
<INTEGER, "42", Line: 1, Col: 9>
<SEMICOLON, ";", Line: 1, Col: 11>
        \end{verbatim} 
        & 
        \begin{verbatim}
<IDENTIFIER, "Count", Line: 1, Col: 1>
[Error] Line: 1, Col: 7, Lexeme: "="
<INTEGER, "42", Line: 1, Col: 9>
<SEMICOLON, ";", Line: 1, Col: 11>
        \end{verbatim} \\
        \hline
    \end{tabular}
    \caption{Output Comparison for Test 1}
\end{table}

\section{Performance Analysis}
\begin{itemize}
    \item \textbf{Manual Scanner}: Incurs overhead from method calls (\texttt{peek}, \texttt{advance}) and object creation during the scanning loop. Complexity is roughly $O(N)$.
    \item \textbf{JFlex Scanner}: Uses a pre-compiled transition table (DFA). This typically offers better raw performance ($O(N)$ with smaller constant factors) as state transitions are array lookups rather than conditional logic.
\end{itemize}

\section{Conclusion}
Both scanners successfully meet the language specifications. The \textbf{JFlex} implementation is stricter and more declarative, making it easier to maintain complex regex rules. The \textbf{Manual} implementation offers finer-grained control over execution flow and error recovery but requires more verbose code.

\end{document}
